% Part: turing-machines
% Chapter: undecidability
% Section: representing-tms

\documentclass[../../../include/open-logic-section]{subfiles}

\begin{document}

\olfileid[pt]{tur}{und}{rep}
\olsection{Representando Máquinas de Turing}

\begin{explain}
A fim de representar máquinas de Turing e seu comportamento por meio de
!!a{sentence} da lógica de primeira ordem, temos de definir uma
linguagem adequada. A linguagem consiste em duas partes: !!{predicate}s
para descrever configurações da máquina, e expressões para numerar os
passos de execução (``momentos'') e as posições na fita.

Introduzimos dois tipos de !!{predicate}s, ambos de 2 lugares: para
cada estado~$q$, !!a{predicate}~$\Obj Q_q$, e para cada símbolo de
fita~$\sigma$, !!a{predicate}~$\Obj S_\sigma$. Os primeiros
permitem-nos descrever o estado de~$M$ e a posição de seu cabeçote de
fita; os segundos permitem-nos descrever o conteúdo da fita.

A fim de exprimir as posições do cabeçote de fita e o número de passos
executados, precisamos de uma maneira de exprimir números. Isso é
feito usando !!a{constant}~$\Obj 0$ e uma função de $1$ lugar~$\prime$,
a função sucessor. Por convenção, ela é escrita \emph{depois} de seu
argumento (e omitimos os parênteses).

Para cada número $n$ há um termo canônico~$\num{n}$, o \emph{numeral}
de~$n$, que o representa em~$\Lang L_M$. $\num{0}$ é $\Obj 0$, $\num{1}$
é $\Obj 0'$, $\num{2}$ é $\Obj 0''$, e assim por diante. Mais
formalmente:
\begin{align*}
\num{0} & = \Obj 0 \\
\num{n+1} &= \num{n}'
\end{align*}
O termo $\num{0}$, isto é, $\Obj 0$, nomeia a posição mais à esquerda
na fita, bem como o instante anterior ao primeiro passo de execução (a
configuração inicial). O termo $\num{1}$, isto é, $\Obj 0'$, nomeia o
quadrado à direita do quadrado mais à esquerda, e o instante posterior
ao primeiro passo de execução, e assim por diante.

Também introduzimos !!a{predicate}~$<$ para exprimir tanto a ordenação
das posições na fita (quando significa ``à esquerda de'') quanto a dos
passos de execução (caso em que significa ``antes'').

Uma vez que tenhamos a linguagem montada, listamos os ``axiomas'' de
$!T(M, w)$, isto é, as !!{sentence}s que, tomadas em conjunto,
descrevem o comportamento de~$M$ quando executada na entrada~$w$.
Haverá !!{sentence}s que estabelecem condições sobre $\Obj 0$,
$\prime$ e $<$, !!{sentence}s que descrevem a configuração de entrada,
e !!{sentence}s que descrevem qual é a configuração de $M$ após
executar uma instrução particular.
\end{explain}

\begin{defn}
  \ollabel{defn:tm-descr}
Dada uma máquina de Turing $M = \tuple{Q, \Sigma, q_0, \delta}$, a
linguagem~$\Lang L_M$ consiste em:
\begin{enumerate}
\item Um !!{predicate} de dois lugares $\Obj Q_q(x, y)$ para todo
  estado~$q \in Q$.  Intuitivamente, $\Obj Q_q(\num{m}, \num{n})$
  exprime ``após $n$ passos, $M$ está no estado~$q$ examinando o
  $m$-ésimo quadrado.''
\item Um !!{predicate} de dois lugares $\Obj S_\sigma(x, y)$ para todo
  símbolo~$\sigma\in \Sigma$.  Intuitivamente, $\Obj S_\sigma(\num{m},
  \num{n})$ exprime ``após $n$ passos, o $m$-ésimo quadrado contém o
  símbolo~$\sigma$.''
\item Um !!{constant} $\Obj 0$
\item Um !!{function} de um lugar $\prime$
\item Um !!{predicate} de dois lugares $<$
\end{enumerate}
\end{defn}

As !!{sentence}s que descrevem a operação da máquina de Turing~$M$ na
entrada $w = \sigma_{i_1}\dots\sigma_{i_k}$ são as seguintes:
\begin{enumerate}
\item Axiomas que descrevem números e~$<$:
\begin{enumerate}
\item !!^a{sentence} que diz que todo número é menor que seu sucessor:
\[
\lforall[x][x < x']
\]
\item !!^a{sentence} que garante que $<$ é transitivo:
\[
\lforall[x][\lforall[y][\lforall[z][
      ((x < y \land y < z) \lif x < z)]]]
\]
\end{enumerate}
\item Axiomas que descrevem a configuração de entrada:
\begin{enumerate}
\item Após $0$~passos---antes de a máquina começar---$M$ está no
  estado inicial~$q_0$, examinando o quadrado~$1$:
\[
\Obj Q_{q_0}(\num{1}, \num{0})
\]
\item Os primeiros $k+1$ quadrados contêm os símbolos $\TMendtape$,
  $\sigma_{i_1}$, \dots, $\sigma_{i_k}$:
\[
\Obj S_\TMendtape(\num{0}, \num{0}) \land
\Obj S_{\sigma_{i_1}}(\num{1}, \num{0}) \land
\dots \land
\Obj S_{\sigma_{i_k}}(\num{k}, \num{0})
\]
\item Caso contrário, a fita está vazia:
\[
\lforall[x][(\num{k} < x \lif \Obj S_\TMblank(x, \num{0}))]
\]
\end{enumerate}
\item Axiomas que descrevem a transição de uma configuração para a
  seguinte:

Para o que segue, seja $!A(x, y)$ a conjunção de todas as
!!{sentence}s da forma
\[
\lforall[z][
  (((z < x \lor x < z) \land \Obj S_\sigma(z, y))
  \lif \Obj S_\sigma(z, y'))]
\]
em que $\sigma \in \Sigma$.  Usamos $!A(\num{m},\num{n})$ para exprimir
``exceto no quadrado~$m$, a fita após $n+1$ passos é a mesma que após
$n$ passos.''
\begin{enumerate}
\item \ollabel{rep-right} Para toda instrução $\delta(q_i, \sigma) =
  \tuple{q_j, \sigma', \TMright}$, a !!{sentence}:
\begin{align*}
& \lforall[x][\lforall[y][(
   (\Obj Q_{q_i}(x, y) \land \Obj S_{\sigma}(x, y)) \lif {}]] \\
&\qquad   (\Obj Q_{q_j}(x', y') \land \Obj S_{\sigma'}(x, y') \land
!A(x, y)))
\end{align*}
Isto diz que se, após~$y$ passos, a máquina está no estado~$q_i$
examinando o quadrado~$x$ que contém o símbolo~$\sigma$, então após
$y+1$ passos ela está examinando o quadrado~$x+1$, está no
estado~$q_j$, o quadrado~$x$ agora contém~$\sigma'$, e todo quadrado
distinto de~$x$ contém o mesmo símbolo que continha após~$y$ passos.

\item \ollabel{rep-left} Para toda instrução $\delta(q_i, \sigma) =
  \tuple{q_j, \sigma', \TMleft}$, a !!{sentence}:
\begin{align*}
& \lforall[x][\lforall[y][
    ((\Obj Q_{q_i}(x', y) \land \Obj S_{\sigma}(x', y)) \lif {}]]\\
& \qquad   (\Obj Q_{q_j}(x, y') \land \Obj S_{\sigma'}(x', y') \land
!A(x, y))) \land {}\\
& \lforall[y][((\Obj Q_{q_i}(\num{0}, y) \land \Obj S_{\sigma}(\num{0},
    y)) \lif {}]\\
& \qquad (\Obj Q_{q_j}(\num{0}, y') \land \Obj S_{\sigma'}(\num{0},
  y') \land !A(\num{0}, y)))
\end{align*}
Reserve um momento para pensar em como isso funciona: agora não
começamos com ``se examinando o quadrado~$x$ \dots'', mas: ``se
examinando o quadrado $x+1$ \dots''. Um movimento para a esquerda
significa que no próximo passo a máquina está examinando o
quadrado~$x$.  Mas o quadrado em que se escreve é~$x+1$.  Fazemos
assim porque não temos subtração nem uma função predecessor.

Note que os números da forma $x+1$ são $1$, $2$, \dots, isto é, isso
não cobre o caso em que a máquina está examinando o quadrado~$0$ e
deveria mover-se para a esquerda (o que, é claro, não pode fazer---ela
simplesmente permanece onde está). Esse caso especial é coberto pela
segunda conjunção: ela diz que se, após $y$ passos, a máquina está
examinando o quadrado~$0$ no estado $q_i$ e o quadrado~$0$ contém o
símbolo~$\sigma$, então após $y+1$ passos ela ainda está examinando o
quadrado~$0$, está agora no estado~$q_j$, o símbolo no quadrado~$0$ é
$\sigma'$, e os quadrados distintos do quadrado~$0$ contêm os mesmos
símbolos que continham após $y$~passos.
\item \ollabel{rep-stay} Para toda instrução $\delta(q_i, \sigma) =
  \tuple{q_j, \sigma', \TMstay}$, a !!{sentence}:
\begin{align*}
& \lforall[x][\lforall[y][(
   (\Obj Q_{q_i}(x, y) \land \Obj S_{\sigma}(x, y)) \lif {}]] \\
&\qquad   (\Obj Q_{q_j}(x, y') \land \Obj S_{\sigma'}(x, y') \land
!A(x, y)))
\end{align*}
\end{enumerate}
\end{enumerate}
Seja $!T(M, w)$ a conjunção de todas as !!{sentence}s acima para a
máquina de Turing~$M$ e a entrada~$w$.

A fim de exprimir que~$M$ eventualmente para, temos de encontrar uma
!!{sentence} que diga ``após algum número de passos, a função de
transição será indefinida.''  Seja $X$~o conjunto de todos os pares
$\tuple{q, \sigma}$ tais que~$\delta(q, \sigma)$ é indefinida.  Seja
então $!E(M, w)$ a !!{sentence}
\[
\lexists[x][\lexists[y][(\bigvee_{\tuple{q, \sigma} \in
      X}(\Obj Q_q(x, y) \land \Obj S_\sigma(x, y)))]]
\]

Se usarmos uma máquina de Turing com um estado de parada
designado~$h$, é ainda mais fácil: então a !!{sentence}~$!E(M, w)$
\[
\lexists[x][\lexists[y][\Obj Q_h(x, y)]]
\]
exprime que a máquina eventualmente para.

\begin{prop}
\ollabel{prop:mlessk}
Se $m < k$, então $!T(M, w) \Entails \num{m} < \num{k}$
\end{prop}

\begin{proof}
Exercício.
\end{proof}

\begin{prob}
Prove \olref[tur][und][rep]{prop:mlessk}.
(Dica: use indução em $k-m$).
\end{prob}

\end{document}
